\section{Reproducibility information}
\label{app:reproducibility}

\subsection{Retained inputs and outputs}

The repository retains the complete Experiment~4 workflow, the 200 accepted
synthetic data sets supplied to both fitting models, combined task-level
results, sampled transmission-rate paths, comparison figures, and selected-task
convergence diagnostics.  The report-specific copy of the R source is stored
under \path{report/code/experiment4_source/}; the canonical executable workflow
is under
\path{experiments/experiment_4_nmif600_model_comparison/code/}.  Numerical
claims in Section~\ref{sec:results} were checked against the retained comparison
tables rather than inferred from the figures.

Each shared-data directory records the observed and simulated series,
simulation metadata, and file checksums.  The manifest at
\path{experiments/experiment_4_nmif600_model_comparison/shared_data/MANIFEST.csv}
provides a compact task index.  The paired comparison code verifies matching
simulation seeds and observed-data checksums before calculating recovery
metrics.

\subsection{Trajectory and likelihood semantics}

Five independent post-IF2 particle filters evaluate each starting-value
endpoint.  Their likelihood estimates are combined using
\texttt{pomp::logmeanexp}, and the endpoint with the largest finite summary is
selected.  This likelihood-scale mean is the only averaging step relevant to
best-fit selection; it is unrelated to averaging latent states.

For the Gamma-noise model, a final 50,000-particle filter retains ancestry and
\texttt{filter\_traj()} extracts one sampled latent path.  The saved
\texttt{B\_path.csv} values at the 70 observation times are used directly in
RSS, RMSE, signed mean error, and absolute overall bias.  No filtering mean
enters those metrics.  The constant-model path is its selected static estimate
repeated at the same times.

\subsection{Exact-reproduction boundary}

The workflow assigns task-specific simulation, IF2, evaluation, and final-filter
seeds.  Nevertheless, the original high-performance-computing runs did not
retain a complete software lockfile.  Exact Monte Carlo reproduction can
therefore depend on the R, \texttt{pomp}, compiler, and random-number-generator
versions.  Recalculation should preserve the supplied seeds, configuration, and
accepted-data checksums, and should report any resulting numerical differences
rather than silently replacing the retained canonical outputs.
